\section{Asymmetric purity and discrepancy}
\label{sec:asymmetric-pure-patches-under-2}

The initial discrepancy estimate \eqref{eq:source-2} will now exclude pure patches whose two widths can be very different. We work along a sequence satisfying that estimate, and all patch laws use the retained labels. Set
\[
 \eta=\min(\eta_0/2,.04),\qquad \tau=\eta/4.
\]
\begin{lemma}[Asymmetric purity exclusion]\label{lem:asymmetric-purity}
Suppose there is a finite balanced tag mixture \(\Lambda\) in one color with laws \(\mu_i\) on \(X\), \(\nu_i\) on \(Y\), of respective widths \(\le n^\gamma,n^\beta\), where \(\gamma<1,\ 0<\beta<\tau/4\) are fixed. Every label in the second support has degree into \(\mu_i\) at least \(1-\epsilon\) with \(\epsilon=\exp(-n^{p})\), \(p>0\) fixed. These hypotheses imply a monochromatic cube. The same works with sides reversed. Balance here means pointwise tag averages \(\le K/N\) as before.
\end{lemma}

\begin{proof}
We construct probability rows for the odd vertices and posterior rows for the even vertices. The required edges will follow from their supports. To obtain injective assignments, we must also keep both sets of column sums bounded.

\noindent\textbf{Step 1: the key grid and trimming.}\quad
Take \(s=\lceil n^\tau\rceil\) chunks of length \(\lfloor n^{.2}\rfloor\), using \(m\) bits altogether. Bin each chunk count into consecutive intervals of mass at most \(2n^{-.04}\), as in the preceding grid construction. Write \(g\) for the vector of bins and \(E(g)\) for its distinct one-bin-step neighbors, so \(|E(g)|\le2s\). Keep the full residual word \(a\in Q_{n-m}\). We use both an even and an odd task at every cell \((g,a)\). The padded even neighbors of an odd task are
\[
 (g,b)\quad\bigl(d_H(a,b)\le1\bigr),
 \qquad (u,a)\quad\bigl(u\in E(g)\bigr).
\]
Call the first set ordinary neighbors and the second set cross neighbors. These lists contain every actual cube edge. For locality statements, \(B_{\rm grid}(g,k)\) denotes the radius-\(k\) ball in the grid graph, and \(B_Q(a,R)\) a residual Hamming ball.

Fix \(\delta=10^{-4}\). We will choose a sufficiently large integer \(h\), independent of \(n\); the list counts below have exponents independent of \(h\). We first trim the mixture so that the survival probability of a first-side label under one independent \(h\)-tuple is nearly constant:
\[
 \alpha_x:=\mathbb E_\Lambda d_G(x,\nu_i)^h
       =2^{-h}(1+o(1/s))
 \quad\text{for every retained first-side label }x.
\]
To see this, let
\[
 b(x)=\Lambda\{i:|d_G(x,\nu_i)-1/2|>2n^{-\eta}\}.
\]
For a fixed \(i\), \eqref{eq:source-2}, applied to uniform laws on either exceptional set, bounds their total size by \(2Ne^{-n^\eta}\). Hence
\[
 \sum_x b(x)\le2Ne^{-n^\eta},\qquad
 D:=\{x:b(x)>e^{-n^\eta/2}\},\qquad
 |D|\le2Ne^{-n^\eta/2}.
\]
Balance gives \(\mathbb E_\Lambda\mu_i(D)\le2K e^{-n^\eta/2}\). Discard tags with \(\mu_i(D)>1/2\), whose total mass is at most \(4K e^{-n^\eta/2}\), normalize the remaining tag law, and replace each \(\mu_i\) by its restriction to \(D^c\). Continue to use the same symbols for these laws. Balance changes only by constant factors, first-side widths increase by \(O(1)\), and each own-color defect is at most \(2\epsilon\). The remaining tag law differs from the original one by \(e^{-\Omega(n^\eta)}\) in total variation. Thus, uniformly over retained \(x\),
\[
 \alpha_x
 =2^{-h}\bigl(1+O(hn^{-\eta})
                  +O(2^h e^{-\Omega(n^\eta)})\bigr)
 =2^{-h}(1+o(1/s)),
\]
because \(h\) is fixed and \(\tau<\eta\).

\noindent\textbf{Step 2: hidden tuples, center laws, and base gates.}\quad
At each grid key sample an independent hidden tuple
\[
 \Theta_g\sim R':=\mathbb E_\Lambda\nu_i^{\otimes h}.
\]
Let \(\eta_g\) be the posterior law of its mixture tag. For the following quantities the key \(g\) is fixed. Define two unnormalized masses under \(\mu_i\):
\[
 \begin{aligned}
 d_i^-&=\mu_i\{x:x\text{ hits every coordinate of }\Theta_u,\
                                     u\in E(g)\},\\
 d_i^+&=\mu_i\{x:x\text{ hits every coordinate of }\Theta_u,\
                                     u\in E(g)\cup\{g\}\},\\
 A_g&=2^{-h|E(g)|},\qquad \Delta=e^{-n^{p/2}}.
 \end{aligned}
\]
Here and below a hit means adjacency in the chosen color. Tilt the tag law by the mass of cross hits, retaining tags for which adding the own tuple costs at most a \(\Delta\)-fraction:
\[
 S_g(i)=\frac{\eta_g(i)d_i^-}{Z_g}
       1[d_i^-\ge e^{-n^{2\tau}},\
                         d_i^+\ge(1-\Delta)d_i^-].
\]
The normalizer is \(Z_g\). These laws are used where their denominators are positive; local fallbacks will be specified below. For a retained tag, the corresponding single-anchor law is
\[
 U_{g,i}(x)=
 \frac{\mu_i(x)1[x\text{ hits }\Theta_u,\
                         u\in E(g)\cup\{g\}]}{d_i^+}.
\]

We next establish the bounds on hidden tuples that will be used throughout this construction. At a fixed key their combined exception probability is at most \(e^{-n^{c'}}\), for some fixed \(c'>0\):
\[
 \eta_g\le e^{h n^\beta+n^{\tau/2}}\Lambda,
 \qquad .8A_g\le Z_g\le1.2A_g.
\]
We also require each of \(\int d_i^-\,d\eta_g(i)\) and \(\int d_i^-\,d\Lambda(i)\) to be within a factor \(1.1\) of \(A_g\). With any one cross key omitted, the analogous requirement is within a factor \(1.1\) of \(2^hA_g\). We call these the base gates.

For the posterior bound, the event
\(R'(\Theta_g)<N^{-h}e^{-n^{\tau/2}}\) has probability at most \(e^{-n^{\tau/2}}\). Outside it,
\[
 \frac{\eta_g(i)}{\Lambda(i)}
 =\frac{\nu_i^{\otimes h}(\Theta_g)}{R'(\Theta_g)}
 \le e^{h n^\beta+n^{\tau/2}}.
\]
The aggregate first laws
\(\overline\mu_\eta=\int\mu_i\,d\eta_g(i)\) and
\(\overline\mu_\Lambda=\int\mu_i\,d\Lambda(i)\) consequently satisfy
\[
 N\max\overline\mu_\eta\le K e^{h n^\beta+n^{\tau/2}},
 \qquad N\max\overline\mu_\Lambda\le K.
\]
Their widths are therefore less than \(n^\eta/2\) for large \(n\), even though individual \(\mu_i\)'s may have much greater width.

Filter either aggregate successively by the independent cross tuples. Conditional on a cross tuple's latent tag, its coordinates are independent with law \(\nu_i\). On preceding regular hits, the aggregate width has increased by only \(O(hs)\), and remains less than \(n^\eta\). A next hit fraction outside \(1/2\pm2n^{-\eta}\) has probability at most \(2e^{-n^\eta}\): restricting the next \(\nu_i\) to a larger exceptional mass would produce, together with the current aggregate, a pair contradicting \eqref{eq:source-2}. Thus
\[
 \prod_{\text{cross coordinates}}
       \left(\tfrac12\pm2n^{-\eta}\right)
 =A_g(1+o(1)).
\]
The same calculation omitting one cross key gives \(2^hA_g(1+o(1))\). A union over both aggregates and all these deletions costs at most \(O(hs^2)e^{-n^\eta}\).

It remains to control the cutoffs in \(S_g\). Put
\(L_g=\int(d_i^--d_i^+)\,d\eta_g(i)\).
Given \(\Theta_g\), every tag in the support of \(\eta_g\) has own-miss mass at most \(2h\epsilon\). A fixed retained \(x\) survives the independent cross tuples with probability
\(\alpha_x^{|E(g)|}=A_g(1+o(1))\). Hence
\[
 \mathbb E[L_g\mid\Theta_g]\le O(h\epsilon A_g).
\]
The weight removed by \(d_i^+<(1-\Delta)d_i^-\) is at most \(L_g/\Delta\), so
\[
 \Pr(L_g/\Delta>.1A_g\mid\Theta_g)
 \le O(h\epsilon/\Delta)
 =\exp(-n^p+n^{p/2}+O_h(1)).
\]
The other cutoff removes at most
\(e^{-n^{2\tau}}=o(A_g)\), since \(\log A_g^{-1}=O(hs)\).
Together with the cross-normalizer bounds these give
\(.8A_g\le Z_g\le1.2A_g\). This calculation uses no ordering between \(p\) and \(\tau\).

\noindent\textbf{Step 3: local centers and the future anchors.}\quad
In each grid slice use the height device on \(Q_{n-m}\), independently sampling prospective positions and activations across slices. Each prospective center in slice \(g\) receives an independent tag with law \(S_g\), conditional on the hidden tuples. Take step length \(D=2\), radius \(r=\lfloor\rho n\rfloor\), and \(\lambda=n^{10}\). The choices
\[
 b=\tau/16,\qquad b_0=b/4,\qquad \rho=.01
\]
are suitable; the height parameters \(\zeta,\sigma,1-\theta\) can be chosen sufficiently small in terms of \(b_0\) to meet Lemma~\ref{lem:height-selection}. At good Lipschitz heights, the ordinary neighbors of an odd cell occupy at most two levels. The crowd bound at one cell of each level then limits their distinct center IDs to \(2n^b\), which is at most
\(T=\lceil n^{\tau/8}\rceil\) for large \(n\).

We specify eligibility below before drawing activations or actual anchors. Once a center has been selected at every even cell, write \(i(g,a)\) for its tag and draw
\(W_{g,a}\sim U_{g,i(g,a)}\), independently between cells at this stage. Cells choosing the same center share its tag, but receive independent anchor samples. This distinction permits the subsequent product likelihoods at odd cells.

\subsection{Presentations at odd cells}
\label{sec:presentations-at-odd-cells}

\noindent\textbf{Step 4: a fixed-presentation reference experiment.}\quad
Consider an odd cell \((g,a)\) with a prescribed list of at most \(T\) distinct own-slice IDs and one ID in each cross slice. The observations are the tags at the own-slice IDs and the pairs \((i,W_{u,a})\) at the cross IDs. Ordinary anchors are not observed. For now these are independent hypothetical observations: conditional on hidden tuples, their laws are \(S_g\) internally and \(S_u(i)U_{u,i}\) externally.

Replace the target tuple \(\Theta_g\) by a candidate \(\xi\) with prior \(R'\), holding all other hidden tuples fixed. Retain the base gates at \(g\) and its grid neighbors; these gates depend on the candidate but not on the hypothetical tag and anchor outcomes. Use the following candidate-independent probability references for the data:
\begin{itemize}
\item For an internal tag, use \(\Lambda(i)d_i^-\), normalized. Here the cross-hit mass \(d_i^-\) does not involve \(\xi\). The density of \(S_g\) against this reference is at most
\(\exp(hn^\beta+n^{\tau/2}+O(1))\).
\item For a cross pair at \(u\in E(g)\), use
\(\eta_u(i)\mu_i(x)\) restricted to hits of the cross tuples of \(u\) other than \(g\), and normalize. Its normalizer is at most \(1.1\cdot2^hA_u\). On the base gate,
\[
 S_u(i)U_{u,i}(x)
 \le\frac{\eta_u(i)\mu_i(x)}
           {(1-\Delta)Z_u}
     1[x\text{ hits }\Theta_w,\ w\in E(u)\cup\{u\}],
\]
because \(d_i^-/d_i^+\le(1-\Delta)^{-1}\). Dividing by the indicated reference therefore costs at most \(2^{h+2}\).
\end{itemize}
The second calculation is the reason for tilting center tags by \(d_i^-\). At undefined reference normalizers use fixed fallbacks independent of the omitted tuple; validity will require the legitimate positive cases.

Let \(Q\) be the product of these references, and let \(F_\xi\) be the product observation density against \(Q\), multiplied by the candidate's base-gate indicator. Define
\[
 M=\int F_\xi\,dR'(\xi),\qquad
 \varepsilon_0=e^{-\delta h s\log n}.
\]
The true-gated data have subdensity \(M\) against \(Q\), so the joint probability of \(M<\varepsilon_0\), including \(M=0\), is at most \(\varepsilon_0\). Require \(M\ge\varepsilon_0\) on a valid presentation. The posterior \(F_\xi\,dR'(\xi)/M\) then has log density against the uniform law on \(Y^h\) at most
\[
 hn^\beta+T(hn^\beta+n^{\tau/2}+O(1))
          +O(hs)+\delta hs\log n
 \le1.5\delta hs\log n.
\]
Here \(T=n^{\tau/8+o(1)}\), \(\beta<\tau/4\), and \(h\) is fixed.

Every candidate in this posterior's support hits all observed cross anchors, coordinate by coordinate. Each observed ordinary tag \(i\) has a positive \(S_g(i)\) factor in \(F_\xi\). Its defining cutoff \(d_i^+\ge(1-\Delta)d_i^-\) therefore gives
\[
 \mu_i\{\text{hits }\xi\mid
                  \text{hits }\Theta_{E(g)}\}\ge1-\Delta.
\]
The conditioning set on the left does not change as \(\xi\) varies. These two support properties will provide the required edges.

\noindent\textbf{Step 5: hidden-history conditioning and local selection.}\quad
We now arrange that the fixed-list denominator test remains likely after selecting centers. There are two successive conditional probabilities. For fixed hidden tuples, let \(q_{g,k}\) be the true-gated probability of \(M<\varepsilon_0\) over the hypothetical tags and cross anchors, with \(k\le T\) distinct internal IDs. The preceding estimate and Markov's inequality give
\[
 \Pr(q_{g,k}>\varepsilon_0^{1/2})
       \le\varepsilon_0^{1/2}.
\]
Only \(k\) matters here, not positions or residual words. Use a local lemma on the independent hidden tuples to impose the base gates and \(q_{g,k}\le\varepsilon_0^{1/2}\) for every \(k\le T\). One combined event per grid key suffices. It uses hidden keys in \(B_{\rm grid}(g,2)\), has probability
\[
 q_H\le C e^{-n^{c'}}+(T+1)\varepsilon_0^{1/2}
        \le e^{-n^{c_H}}
\]
for some \(c_H>0\), and dependency degree \(O((1+2s)^4)\). Thus charges \(x_H=2q_H\) suffice for large \(n\). All posterior formulas continue to use the raw references of Step 4, rather than references conditioned on this avoidance event.

Fix such a hidden history and generate prospective positions and tags. For a candidate ID list at an odd cell, let \(q_L\) be the denominator-failure probability over hypothetical cross anchors after fixing its tags. The hidden-history bound gives
\[
 \Pr_{\rm tags}(q_L>\varepsilon_0^{1/4})
       \le\varepsilon_0^{1/4}.
\]
Call precisely these lists bad. At each odd cell consider all lists using present IDs in the balls of its incident even cells, at all levels, with at most \(T\) internal IDs and one per cross key. No assignment of internal IDs to particular ordinary neighbors is needed. On prospective ball counts between \(\lambda/2\) and \(2\lambda\), their number is at most
\[
 L_n=\exp(C(s+T)\log n),
\]
where \(C\) is independent of \(h\). Lists with disjoint ID sets have independent tag outcomes, conditional on the hidden tuples and positions. Therefore
\[
 \Pr(\text{\(n\) disjoint bad lists at a cell})
 \le L_n^n\varepsilon_0^{n/4}
 =\exp\!\left(n[C(s+T)-\delta hs/4]\log n\right).
\]
Choose the fixed integer \(h\) so that \(\delta h/4>2C\). Since \(T=o(s)\), the exponent is negative of order \(ns\log n\), which dominates the logarithm \(O(n+s\log n)\) of the number of cells.

Choose a maximal disjoint family of bad lists at each odd cell, using a fixed ordering of its local lists, and mark their IDs forbidden at the incident even sites whose corresponding balls contain them. On the event just proved, each family contains fewer than \(n\) lists. An even site has \(O(n+s)\) incident odd cells, so the total number of forbidden IDs per site-level is at most
\[
 O(n(n+s)(s+T))=o(\lambda).
\]
Prospective ball counts lie in the required range simultaneously with high probability by binomial tails. At least \(\lambda/3\) choices therefore remain eligible at every site-level. Any used list whose IDs are all eligible cannot be bad: it would be disjoint from the marked maximal family.

Only now activate centers and select them by the long height rule. Eligibility is independent of every slice's activations. Lemma~\ref{lem:height-selection} gives good Lipschitz heights simultaneously in all slices with probability \(1-o(1)\): its failure exponent exceeds \(n\) by a power, whereas there are only \(\exp(O(s\log n))\) grid keys. In particular every used list passes the conditional bound \(q_L\le\varepsilon_0^{1/4}\), and ordinary fans have at most \(T\) IDs.

We record the inputs of this selection rule. Eligibility at \((t,b)\) examines only the bad-list families at its incident odd cells. Their lists read center positions and tags in
\[
 B_{\rm grid}(t,2)\times B_Q(b,r+2)
\]
and hidden tuples in \(B_{\rm grid}(t,3)\). The conditional list test uses the fixed-presentation experiment, not any selected IDs or adjusted label kernels. The height rule stays within slice \(t\) and consults statuses at residual distance at most \(4H\). Consequently selection at \((t,b)\) uses center positions, tags, activations and ties within
\[
 B_{\rm grid}(t,2)\times B_Q(b,r+4H+2),
 \qquad\text{and hidden tuples in }B_{\rm grid}(t,3).
\]
All levels are included. This is a finite sequence of consultations: a list test never asks for another selection. Define the same local rules off successful histories, using local fallbacks when a law or test is undefined and independently randomized uniform choices among eligible active IDs when usable. Such fallbacks do not inspect success elsewhere.

\noindent\textbf{Step 6: selection adjustment under candidate replacement.}\quad
Fix positions and all hidden tuples except \(\Theta_g\). In a raw experiment draw \(\Theta_g\sim R'\), generate the center tags and selections needed at the incident even cells under their resulting laws, and then draw their cross anchors. There is no further conditioning. A valid presentation requires the true base gate of Step 4, legitimate local selections with at most \(T\) internal IDs, incident-ball position counts at all levels at most \(2\lambda\), and \(M\ge\varepsilon_0\). Ordinary anchor values are still absent.

Group presentations by their distinct internal IDs and their one cross ID at each key; call this list \(L\). Its observed tags and cross pairs have product reference \(Q_L\) from Step 4. Given a candidate \(\xi\), the subevent of validly presenting \(L\) has density
\[
 F_\xi a_\xi\quad\text{against }Q_L,\qquad 0\le a_\xi\le1.
\]
Indeed the listed centers must first have the observed independent tag outcomes. Selection and validity can only reduce their probability. After selection, the cross anchors have the indicated independent laws conditional on the selected tags. Thus the unselected product density \(F_\xi\) dominates the valid-presentation subdensity.

The factor \(a_\xi\) is obtained by disintegrating this raw experiment. In particular, when \(\Theta_g\) changes, every affected center law \(S_u\) changes with it. The unobserved center tags and the local selections are integrated under those new laws, rather than held at their realized values. The cross-key reference deletes the entire dependence on \(\Theta_g\). The different lists are disjoint outcomes of this same candidate-dependent experiment, so
\[
 \sum_L\int F_\xi a_\xi\,dQ_L\le1
 \quad\text{for each }\xi.
\]

\noindent\textbf{Step 7: posterior truncation and mean recovery.}\quad
For the realized presentation put
\(M^a=\int F_\xi a_\xi\,dR'(\xi)\).
Use the selected-presentation posterior
\(F_\xi a_\xi\,dR'/M^a\) when \(M^a\ge\varepsilon_0 M\);
otherwise use the base posterior \(F_\xi\,dR'/M\).
In either case its log density against uniform tuples is bounded by
\[
 hn^\beta+T(hn^\beta+n^{\tau/2}+O(1))
       +O(hs)+2\delta hs\log n
 <.003hs\log n.
\]
The second denominator costs at most one further factor \(\varepsilon_0^{-1}\).

Let \(q\) be the average coordinate marginal of this tuple posterior and
\(\mathcal H=\{y:Nq(y)>e^{.01s\log n}\}\).
Then \(|\mathcal H|/N\le e^{-.01s\log n}\), and the joint density cap implies
\[
 \Pr(\text{at least \(h/2\) coordinates lie in }\mathcal H)
 \le 2^h e^{.003hs\log n}(|\mathcal H|/N)^{h/2}
 =o(1).
\]
It follows that \(q(\mathcal H)\le1/2+o(1)\). Delete \(\mathcal H\) and normalize; the normalization factor is bounded by a constant. Denote this law by \(p^0_{g,a}\), setting it to zero on invalid presentations.

For later loads, the important conclusion is its raw mean bound
\[
 \mathbb E p^0_{g,a}(y)
       \le O(1)\mathbb E_\Lambda\nu_i(y)\le O(K)/N.
\]
Here the expectation is in the experiment of Step 6, at fixed positions and non-target hidden tuples. To verify it, valid data for list \(L\) have subdensity \(M_L^a\) against \(Q_L\). On the selected-posterior cases, cancellation gives the following inequality of measures on the candidate tuple:
\[
 \sum_L\int_{\{M_L^a\ge\varepsilon_0M_L\}}
 M_L^a\,\frac{F_\xi a_\xi\,dR'(\xi)}{M_L^a}\,dQ_L
 \le dR'(\xi)\sum_L\int F_\xi a_\xi\,dQ_L
 \le dR'(\xi).
\]
On the other cases, \(M_L^a<\varepsilon_0M_L\) bounds their contribution by
\[
 \varepsilon_0\,dR'(\xi)\sum_L\int F_\xi\,dQ_L
 \le\varepsilon_0 L_n\,dR'(\xi)=o(dR'(\xi)).
\]
The list count follows from the local position gate, and \(h\) was chosen large enough that \(\varepsilon_0L_n=o(1)\). Taking average coordinate marginals and paying the bounded truncation factor proves the claim.

The input scope of this whole computation will matter when its means are multiplied. The incident even selections are at grid and residual distance at most one from \((g,a)\). By Step 5, their raw center inputs lie in
\[
 B_{\rm grid}(g,3)\times B_Q(a,r+4H+3),
 \qquad\text{and their hidden inputs lie in }B_{\rm grid}(g,4).
\]
The only realized anchors read by \(p^0_{g,a}\) are \(W_{u,a}\), \(u\in E(g)\). Its evaluation of \(a_\xi\) integrates unobserved tags and selection randomness inside the displayed region, at the fixed positions there. Candidate replacement changes these local laws and tests but does not enlarge their scope. Thus, after cross-anchor integration, the composite raw experiment has grid radius \(4\) and residual radius \(r+4H+3=r+o(n)\), and reads no realized anchors. Neither its validity gate nor \(a_\xi\) uses the later anchor-avoidance tests. We use this raw-computed posterior rule on the constrained histories as well.

\noindent\textbf{Step 8: the ordinary-anchor hit test.}\quad
The support of \(p^0_{g,a}\) already hits every cross anchor. Now restrict it to labels hitting the realized anchors at all padded ordinary even neighbors, and require the retained mass to be at least \(.98\). Call the resulting normalized law \(p_{g,a}\).

To bound this test's failure probability, fix valid presentation data and cross anchors before conditioning the anchors. A support label of \(p^0_{g,a}\) belongs to a coordinate of a supported candidate tuple. For any observed internal tag \(i\), Step 4 bounds its miss probability under \(\mu_i\) conditioned on cross hits by \(\Delta\). Conditioning additionally on the true own tuple increases this by at most \((1-\Delta)^{-1}\). Thus its miss probability under \(U_{g,i}\) is at most \(\Delta/(1-\Delta)\). Summing over the \(O(n)\) ordinary anchors and applying Markov's inequality gives failure probability \(O(n\Delta)\). On validity,
\[
 p_{g,a}\le p^0_{g,a}/.98,\qquad
 N\max p_{g,a}\le e^{.02s\log n},
\]
and this law hits every incident actual even anchor. The reference \(p^0_{g,a}\) still omits all ordinary anchor values.

\subsection{Taking the assignment}
\label{sec:taking-the-assignment-asymmetric}

\noindent\textbf{Step 9: selected-anchor load bounds.}\quad
Before sampling actual anchors we need, with high probability on successful hidden and selection histories,
\[
 \operatorname*{avg}_{v\in A}
       N U_{g(v),i(g(v),a(v))}(x)=O(1)
 \quad\text{for every }x.
\]
There are two successive averages to control: over center selection at fixed hidden tuples, and over the hidden tuples themselves.

For the first, retain local validity, including the eligibility-size bounds on the height consultation domain. Given hidden tuples, force a specified center present and fix its tag. The height lemma bounds its gated selection probability, averaged over positions, by \(O(1/\lambda)\) at level zero and \(e^{-n^c}\) at a positive level. The sum over possible centers has weight
\[
 \lambda O(1/\lambda)+H\lambda e^{-n^c}=O(1).
\]
On the base gate the tag and anchor laws satisfy
\[
 S_g(i)U_{g,i}(x)
 \le
 \frac{\eta_g(i)\mu_i(x)\,
       1[x\text{ hits }\Theta_{E(g)}]}
      {(1-\Delta)Z_g}.
\]
Choose a fixed constant \(C_1\) large enough for the preceding selection bound, and define, on the base gate,
\[
 B_g(x):=
 C_1A_g^{-1}1[x\text{ hits }\Theta_{E(g)}]
                    \int N\mu_i(x)\,d\eta_g(i).
\]
Set \(B_g(x)=0\) off that gate. The conditional mean of the gated normalized selected law is at most \(B_g(x)\), whose cap is
\[
 L_B=\exp(O(hs+hn^\beta+n^{\tau/2})).
\]
Under the raw hidden law, averaging \(\Theta_g\) returns the posterior tag law to its prior, \(\mathbb E\eta_g=\Lambda\). Independently averaging the cross tuples gives
\(\alpha_x^{|E(g)|}=O(A_g)\). Therefore
\[
 \mathbb E_{\rm raw}B_g(x)
 \le O(1)A_g^{-1}\alpha_x^{|E(g)|}
                    \mathbb E_\Lambda N\mu_i(x)
 =O(K).
\]

Apply Lemma~\ref{lem:scattered-moments} first to these comparison means on the hidden-history avoidance law. Declare rows near when their grid keys are within distance eight. Their fraction is at most
\(f_{\rm grid}=n^{O(1)}(2n^{-.04})^s\), and
\(nf_{\rm grid}L_B=o(1)\).
For separated rows the relevant hidden neighborhoods are disjoint. Removing hidden constraints touching one such neighborhood costs
\(\exp(O((1+2s)^3x_H))=1+o(1)\);
the remaining raw expectations are bounded by \(O(K)\). The scattered-moment estimate and a union over \(x\) show that the row average of \(B_g(x)\) is bounded simultaneously for all \(x\), with high probability.

Fix a hidden history with these bounds and now average the selected laws over center randomness. Their cap is
\[
 L_U=\exp(n^\gamma+n^{2\tau}+O(1)),
\]
by the lower cutoff on \(d_i^-\) and \(d_i^+\ge(1-\Delta)d_i^-\).
Two selections have disjoint center inputs when their residual words are more than \(2r+8H+4\) apart, by Step 5. Up to a constant parity factor, the fraction of nearer words is at most
\[
 f_{\rm res}=2^{-(n-m)}
       \sum_{j\le2r+8H+4}\binom{n-m}{j}
       =e^{-\Omega(n)}.
\]
Here \(r=.01n+O(1)\), \(H=o(n)\), and the binomial-sum exponent is strictly less than \(n\log2\). Since \(\log L_U=o(n)\), we have \(nf_{\rm res}L_U=o(1)\). Conditional independence at separated words and the comparison means \(B_g(x)\) now give the asserted selected-anchor load bound by a second application of the scattered-moment lemma. Both applications keep successful histories as indicators until the raw comparisons have been made.

\noindent\textbf{Step 10: anchor avoidance and predictive alarms.}\quad
Fix a successful hidden and selection history with the selected-law load bound. Failures of these earlier stages may be treated as aborts; their laws will not be reweighted by later avoidance probabilities. Start with independent anchors \(W_{g,a}\sim U_{g,i(g,a)}\). At each odd cell the denominator test fails with probability at most \(\varepsilon_0^{1/4}\), by its selected list's bound, and the ordinary-hit test has the \(O(n\Delta)\) bound of Step 8. We will avoid these events together with the following predictive alarms.

At an actual even position \(v\), write \(U_v\) for its selected anchor law and let its anchor \(z\) vary, holding all other anchors fixed. For the vector \(\boldsymbol y\) of labels at its actual odd neighbors set
\[
 F_z(\boldsymbol y)
   =1_{\mathcal V}\prod_{u\sim v}p_{g(u),a(u)}(y_u),
 \qquad
 M_v(\boldsymbol y)=\int F_z(\boldsymbol y)\,dU_v(z).
\]
Here \(\mathcal V\) requires those neighboring presentations and ordinary-hit tests to be valid; it contains no predictive alarms. All these gates and kernels, including the \(p^0\) lookups at observed cross anchors, are reevaluated when \(z\) varies. Center histories remain fixed.

Choose a product reference \(Q_v\) independent of \(z\). At a neighbor of the same grid key use \(p^0\), or a fixed independent fallback if necessary: \(z\) is an ordinary anchor there and hence was omitted in the definition of \(p^0\). At all other neighbors use the uniform law. There are at most \(m\) such cross neighbors, so
\[
 F_z\le .98^{-n}\bigl(e^{.02s\log n}\bigr)^mQ_v
       \le e^{.03n}Q_v,
 \qquad ms\log n=o(n).
\]
Define predictive failure by
\(M_v=0\) or \(M_v<e^{-.04n}Q_v\).
The joint probability of this failure on \(\mathcal V\), under the raw target anchor and product neighbor sampling, is at most \(e^{-.04n}\), since the data subdensity is \(M_v\). Markov's inequality therefore bounds by \(e^{-.02n}\) the probability over anchors that its conditional product-sampling probability exceeds \(e^{-.02n}\).

At a fixed cell only \(\exp(O(s\log n))\) actual multiplicity profiles occur: the numbers of special-bit neighbors in each cross bin or the own bin determine the formula up to a permutation of the sampled coordinates. Combine their alarm failures into one event per cell. Its probability is at most \(\exp(O(s\log n)-.02n)\). In the cell graph with grid or residual one-step moves, denominator and hit events have anchor radius one and predictive events have anchor radius two. The maximum degree of this graph is \(O(n)\), so the grouped events have dependency degree \(O(n^4)\). A uniform event bound is
\[
 q_A\le \varepsilon_0^{1/4}+O(n\Delta)
                +\exp(O(s\log n)-.02n)
       \le e^{-n^{c_A}}
\]
for some \(c_A>0\). The conditional local lemma applies with charges \(x_A=2q_A\), at each fixed entering history. Sample the anchors from its avoidance law.

\noindent\textbf{Step 11: odd loads by successive removal of constraints.}\quad
We show that the odd kernel column sums are at most \(\theta_0\), as required by the clock lemma, with high probability on success. Let \(\mathcal S\) denote successful earlier histories and selections, and consider \(l\le n\) odd rows \(b_j=(g_j,a_j)\) with pairwise grid distance greater than eight.

First use \(p_{b_j}\le p^0_{b_j}/.98\). One anchor variable is touched by \(O(n^2)\) of the grouped events from Step 10. Removing all events touching the at most \(2s\) cross anchors of a row costs
\(\exp(O(sn^2x_A))=1+o(1)\).
The cross anchors can then be integrated independently under their selected raw laws. No ordinary anchor remains in the integrand.

Next discard global selection success while retaining each local valid-presentation indicator in \(p^0\). Finally remove the hidden-history constraints touching the target tuples \(\Theta_{g_j}\). At most \(O((1+2s)^2)\) hidden events touch one target, so this costs
\(\exp(O((1+2s)^2x_H))=1+o(1)\) per row.
The radius-four scopes of Step 7 are disjoint. Conditional on positions and the other hidden tuples, the target tuples have their independent priors and the local center experiments use disjoint primitive inputs. Each experiment regenerates its center laws and selections when its candidate tuple changes, exactly as in Step 6. The raw mean bound of Step 7 therefore gives
\[
 \mathbb E\!\left[1_{\mathcal S}
                  \prod_{j=1}^l Np_{b_j}(y)\right]
 \le
 \left(\frac{O(K)}{.98}
  \exp\{O(sn^2x_A+(1+2s)^2x_H)\}\right)^l
 \le C^l.
\]
Both charge sums are \(o(1)\), even after multiplication by \(n\).

For the remaining, grid-near rows, the cap is
\(L=e^{.02s\log n}\) and
\[
 nf_{\rm grid}L
 \le n^{O(1)}(2n^{-.04})^s e^{.02s\log n}
 =\exp(-.02s\log n+O(s+\log n))=o(1).
\]
Lemma~\ref{lem:scattered-moments}, followed by Markov and a union over labels, now bounds every normalized average odd load by a constant with high probability. The actual column sums are \(O(2^n/N)=o(1)\), because \(C_n\to\infty\), and are therefore at most \(\theta_0\). This is a joint estimate with earlier successes, rather than a claim that conditioning on those successes preserves independence.

\noindent\textbf{Step 12: injection, even loads, and Hall.}\quad
On a successful prehistory with the odd column bound, apply Lemma~\ref{lem:clock-sampling} with \(B=3\) to the actual odd rows. The predictive-failure predicate at an even position involves its at most \(n\) odd neighbors, each odd row occurs in at most \(n\) such predicates, and their product-law failure probabilities are at most \(e^{-.02n}\). The atom cap from Step 8 is more than sufficient for the lemma's polynomial requirement. We obtain an odd injection avoiding all predictive failures and retaining the joint upper comparison on at most \(n^2\) queries.

At an even position use the posterior row
\[
 w_v(x;\boldsymbol y)
       =\frac{U_v(x)F_x(\boldsymbol y)}{M_v(\boldsymbol y)}.
\]
Its support consists of common neighbors of the assigned odd labels. In auxiliary experiments set it to zero unless local validity and predictive success hold. Since
\[
 N\max U_v\le e^{n^\gamma+n^{2\tau}+O(1)}=e^{o(n)},
\]
the likelihood comparison \(F_x\le e^{.03n}Q_v\) and predictive lower bound \(M_v\ge e^{-.04n}Q_v\) give
\(N\max w_v\le e^{.1n}\) for large \(n\).

For its load estimate fix the entire center history, including the selected-law bound of Step 9; only anchors and odd outputs are varied. Select at most \(n\) even rows whose residual words have pairwise distance greater than four. Their odd-neighbor sets are disjoint, and their anchor calculations have residual radius at most two. Transfer their neighbor outputs jointly to product odd sampling, keeping the successful-prehistory indicator for this comparison. Then discard nonlocal requirements on the anchors and odd outputs and remove the anchor constraints touching each target anchor. Each removal costs \(\exp(O(n^2x_A))=1+o(1)\), and another retained row's calculation does not consult that target.

Fix all other anchors. After integrating the target anchor, the product neighbor data with the local validity indicator have subdensity \(M_v\). Thus posterior cancellation gives, for each label \(x\),
\[
 \sum_{\boldsymbol y}
 M_v(\boldsymbol y)\,
 w_v(x;\boldsymbol y)\,
 1[\text{predictive success}]
 \le U_v(x)\sum_{\boldsymbol y}F_x(\boldsymbol y)
 \le U_v(x).
\]
The same identity multiplies over the separated rows. Their comparison means are thus the selected laws whose average was bounded in Step 9. Residual-near rows have fraction at most \(n^{O(1)}2^{-(n-m)}\), and
\[
 n\cdot n^{O(1)}2^{-(n-m)}e^{.1n}=o(1).
\]
The scattered-moment estimate shows that the joint probability of prior success and any normalized even load exceeding a fixed constant is \(o(1)\). Intersecting its complement with the earlier high-probability events gives a successful history and odd injection for which all even loads are bounded. Their actual column sums are \(O(2^n/N)=o(1)\), hence at most one for large \(n\). Hall's theorem supplies distinct even representatives. Together with the odd injection and the common-neighbor support of every even row, this gives the monochromatic cube.
\end{proof}

\begin{corollary}[Asymmetric power discrepancy]\label{cor:asymmetric-discrepancy}
Consequently in a bad sequence with \eqref{eq:source-2}, after further stabilization for countably many budgets, we can assume:
for every rational \(0<\beta<\tau/4,\ 0<\gamma<1\) there is a constant \(h_{\beta,\gamma}>0\) such that on retained sets all pairs at widths \(\le n^\beta,n^\gamma\) respectively have both-color densities \(1/2\pm n^{-h_{\beta,\gamma}}\) eventually, and similarly in reversed orientation. The exponent may depend on the budget pair; no uniform exponent over all rational pairs is asserted.
\end{corollary}

\begin{proof}
We may replace \(\gamma\) by \(\max(\beta,\gamma)\), so assume \(\beta\le\gamma\). Apply Lemma~\ref{lem:bias-purity} at the budgets \(n^\beta,n^\gamma\), and write \(h_*>0\) for its bias and purity exponent. Choose
\(\beta<\beta^+<\tau/4\) and \(\gamma<\gamma^+<1\).
For large \(n\), the doubled budgets \(2n^\beta,2n^\gamma\), including an \(O(1)\) trimming cost, fit within \(n^{\beta^+},n^{\gamma^+}\). Stabilize the bias property with threshold \(n^{-h_*}\) and the purity property with defect \(e^{-n^{h_*}}\) at the doubled budgets.

If that purity is available, pass to one available color and trim the small-budget side for minimum degree with defect \(e^{-n^{h_*/2}}\). A balanced mixture then meets Lemma~\ref{lem:asymmetric-purity} with the enlarged exponents \(\beta^+,\gamma^+\), placing the small-budget side on its odd side. That lemma gives a cube, so the purity alternative is absent. An available bias of at least \(n^{-h_*}\) at the original budgets would now give a cube by Lemma~\ref{lem:bias-purity}. Stabilization excludes that alternative too. Applying this to each rational budget pair, and to the reversed orientation, gives the stated discrepancy; the exponent is allowed to depend on the pair.
\end{proof}
